\chapter{The Riemann Mapping Theorem}\label{ch:riemannmap}

\begin{center}\itshape
The disc, a square, the inside of a snowflake, the slit plane:\\
to a holomorphic bijection they are all the same.
\end{center}

The preceding chapter carried the upper half-plane onto the unit disc by a map one can write in a line. This chapter proves that map, then Riemann's theorem of 1851: every connected open set in the plane, other than the plane itself, on which holomorphic functions have square roots, admits a holomorphic bijection onto the disc. A square is a map that is not fractional linear. The general domain is found by maximising a derivative, once bounded holomorphic functions are known to have convergent subsequences.

In his 1851 thesis Riemann claimed the theorem for every simply connected region other than the plane. His argument minimised an energy, the ``Dirichlet principle'', and in 1870 Weierstrass showed that such minima need not exist. The first complete proof was Osgood's, in 1900; the one below combines Koebe's square-root trick (1912) with an extremal argument of Fej\'er and Frigyes Riesz, published by Tibor Rad\'o in 1923.

\section{The half-plane, written down}

\begin{lemma}[the Cayley map]\label{lem:cayley}
The map $\varphi(z)=\frac{z-i}{z+i}$ is a holomorphic bijection of the upper half-plane onto the open unit disc, and its inverse $w\mapsto i\frac{1+w}{1-w}$ is holomorphic.
\end{lemma}
\begin{proof}
$|z-i|<|z+i|$ exactly when $z$ is closer to $i$ than to $-i$, that is, when $\operatorname{Im}z>0$, so $\varphi$ maps the upper half-plane into the disc. For the inverse, $\operatorname{Im}\bigl(i\frac{1+w}{1-w}\bigr)=\frac{1-|w|^2}{|1-w|^2}>0$ whenever $|w|<1$, so the inverse formula maps the disc into the upper half-plane. The two maps undo each other by direct expansion.
\end{proof}

A square asks for more than a fractional linear map. The integral $f(z)=\int_0^z(1+w^4)^{-1/2}\,\dd w$ sends the disc onto a square: angles are preserved at every interior point, and the four corners are the images of the four points where $w^4=-1$.

\begin{figure}[ht]
\centering
\fig{ch09-square}
\caption{A Riemann map made explicit: $f(z)=\int_0^z\frac{\dd w}{\sqrt{1+w^4}}$ maps the unit disc onto a square (a Schwarz--Christoffel integral). The images of the circles $|z|=0.2,0.4,\dots,0.99$ and of
sixteen radii, computed from the power series $f(z)=\sum_k\binom{-1/2}k\frac{z^{4k+1}}{4k+1}$. Angles are preserved everywhere inside; at the four corners, images of the points where $w^4=-1$, the right
angles of the square are created from the straight angle of the circle.}
\label{fig:squaremap}
\end{figure}

The general domain has no formula. Three tools find the map: a compactness theorem for bounded holomorphic functions, the Schwarz lemma, and the automorphisms of the disc.

\begin{theorem}[Montel, 1907]\label{thm:montel}
Let $\Omega$ be open and $(f_n)$ holomorphic functions on $\Omega$, uniformly bounded on each compact subset of $\Omega$. Then a subsequence converges uniformly on every compact subset of $\Omega$, to a holomorphic
function.
\end{theorem}
\begin{proof}
Let $K_j=\{z:|z|\le j,\ d(z,\C\setminus\Omega)\ge\frac1j\}$ (with $d(z,\emptyset)=\infty$). These are compact, increase to $\Omega$, and every compact $K\subset\Omega$, which is bounded and at positive distance from
$\C\setminus\Omega$, lies in some $K_j$. For $z\in K_j$, $\overline D(z,\frac1{2j})\subset K_{2j}$. Let $M_j$ bound all $|f_n|$ on $K_{2j}$. If $z,z'\in K_j$ and $|z-z'|<\frac1{4j}$, then every $\zeta\in[z,z']$ has
$\overline D(\zeta,\frac1{4j})\subset\overline D(z,\frac1{2j})\subset K_{2j}$, so $|f_n'(\zeta)|\le4jM_j$ by Cauchy's estimates (Volume~5), and $|f_n(z)-f_n(z')|\le4jM_j|z-z'|$. So $(f_n)$ is bounded and
equicontinuous on $K_j$, and by Arzel\`a--Ascoli (Volume~7) has a subsequence converging uniformly on $K_j$. Extract successively for $K_1,K_2,\dots$ and take the diagonal subsequence: it
converges uniformly on every $K_j$, hence on every compact set. The limit is holomorphic (Volume~5).
\end{proof}

\begin{lemma}[injective holomorphic maps]\label{lem:injective}
Let $f$ be holomorphic and injective on an open set $\Omega$. Then $f'$ has no zeros, $f(\Omega)$ is open, and $f^{-1}:f(\Omega)\to\Omega$ is holomorphic. If $\Omega$ is connected and $f_n$ are injective
holomorphic functions on $\Omega$ converging to $f$ uniformly on compact sets, then $f$ is injective or constant.
\end{lemma}
\begin{proof}
Let $z_0\in\Omega$, $w_0=f(z_0)$, and let $k\ge1$ be the order of the zero of $f-w_0$ at $z_0$ ($f$ is not constant near $z_0$, being injective). Choose $r>0$ with $\overline D(z_0,r)\subset\Omega$, such that $f-w_0$
has no other zero, and $f'$ no zero except possibly $z_0$, in this closed disc (zeros of a non-zero holomorphic function are isolated, Volume~5). Let $\mu=\min_{|z-z_0|=r}|f(z)-w_0|>0$. For $|w-w_0|<\mu$, Rouch\'e's theorem (Volume~5, with $f-w_0$ and the constant $w_0-w$) says that $f-w$ has $k$ zeros in $D(z_0,r)$, with multiplicity. For $w\ne w_0$ they are simple ($f'\ne0$ there), so they are $k$
distinct points, and injectivity forces $k=1$: $f'(z_0)\ne0$. With $k=1$, every $w$ with $|w-w_0|<\mu$ is a value of $f$ on $D(z_0,r)$. This proves that $f(\Omega)$ is open, and also that
$f^{-1}$ is continuous at $w_0$ (as $r$ can be taken arbitrarily small). Then $\frac{f^{-1}(w)-f^{-1}(w_0)}{w-w_0}=\frac{z-z_0}{f(z)-f(z_0)}\to\frac1{f'(z_0)}$ as $w\to w_0$, with $z=f^{-1}(w)\to z_0$.

For the limit statement, suppose $f(z_1)=f(z_2)=w$ with $z_1\ne z_2$ and $f$ not constant. Choose disjoint closed discs around $z_1$ and $z_2$ in $\Omega$ on whose boundary circles $f-w$ has no zero
(identity theorem), and let $\mu>0$ be the minimum of $|f-w|$ on the two circles. For large $n$, $|f_n-f|<\mu$ on both circles, and by Rouch\'e $f_n-w=(f-w)+(f_n-f)$ has a zero in each disc: $f_n$ is not
injective.
\end{proof}

\begin{lemma}[Schwarz; disc automorphisms]\label{lem:schwarz}
Let $D$ be the open unit disc.
\begin{enumerate}
\item[(a)] If $f:D\to D$ is holomorphic with $f(0)=0$, then $|f(z)|\le|z|$ and $|f'(0)|\le1$. If $|f'(0)|=1$, or $|f(z)|=|z|$ for some $z\ne0$, then $f(z)=cz$ with $|c|=1$.
\item[(b)] For $|a|<1$, $\varphi_a(z)=\frac{z-a}{1-\overline az}$ is a holomorphic bijection of $D$ onto $D$, with inverse $\varphi_{-a}$ and $\varphi_a(a)=0$.
\end{enumerate}
\end{lemma}
\begin{proof}
(a) $g(z)=\frac{f(z)}z$ extends holomorphically to $D$ with $g(0)=f'(0)$. On $|z|=r<1$, $|g|\le\frac1r$, so $|g|\le\frac1r$ on $|z|\le r$; let $r\to1$: $|g|\le1$. If $|g|=1$ somewhere in $D$, the maximum principle
makes $g$ constant. (b) $|1-\overline az|^2-|z-a|^2=(1-|a|^2)(1-|z|^2)>0$ for $|z|<1$, so $\varphi_a$ maps $D$ into $D$; and $\varphi_{-a}(\varphi_a(z))=\frac{(z-a)+a(1-\overline az)}{(1-\overline az)+\overline a(z-a)}=\frac{(1-|a|^2)z}{1-|a|^2}=z$, and
symmetrically.
\end{proof}

\begin{theorem}[Riemann mapping theorem]\label{thm:riemannmap}
Let $\Omega\ne\C$ be a connected open set on which every holomorphic function without zeros has a holomorphic square root; Volume~8 proves that this holds whenever $\Omega$ is homologically simply connected. Let $z_0\in\Omega$. Then there is exactly one holomorphic bijection $F:\Omega\to D$ with $F(z_0)=0$ and $F'(z_0)>0$, and $F^{-1}$ is holomorphic.
\end{theorem}
\begin{proof}
Let $\mathcal F$ be the set of injective holomorphic $f:\Omega\to D$ with $f(z_0)=0$.

\emph{$\mathcal F$ is not empty.} Take $a\notin\Omega$ and a holomorphic $q$ with $q(z)^2=z-a$. Then $q$ is injective, and $q(z)\ne-q(w)$ for all $z,w$: equal squares give $z=w$, then $q(z)=0$ and $z=a$. The open
set $q(\Omega)$ (Lemma~\ref{lem:injective}) contains a disc $D(c,r)$, so $q(\Omega)$ misses $D(-c,r)$ and $f_0=\frac r{q+c}$ is injective with $|f_0|\le1$, hence $|f_0|<1$ as $f_0(\Omega)$ is open. Then $\varphi_{f_0(z_0)}\circ f_0\in\mathcal F$.

\emph{An extremal map.} If $\overline D(z_0,\delta)\subset\Omega$, Cauchy's estimate gives $|f'(z_0)|\le\frac1\delta$ for $f\in\mathcal F$. Let $s=\sup_{\mathcal F}|f'(z_0)|$, which is positive by Lemma~\ref{lem:injective}, and $f_n\in\mathcal F$ with
$|f_n'(z_0)|\to s$. By Montel, a subsequence converges uniformly on compact sets to a holomorphic $F$, with $|F'(z_0)|=s$ (a locally uniform limit of holomorphic functions is holomorphic, Volume~5). So $F$ is not constant, hence injective by
Lemma~\ref{lem:injective}; $|F|\le1$ and $F(\Omega)$ is open, so $F(\Omega)\subset D$; and $F(z_0)=0$. So $F\in\mathcal F$ and $|F'(z_0)|$ is maximal.

\emph{$F$ is onto.} Suppose $w\in D\setminus F(\Omega)$. Then $\varphi_w\circ F$ has no zeros, so it has a holomorphic square root $g$, injective with values in $D$. Let $b=g(z_0)$ and $h=\varphi_b\circ g\in\mathcal F$.
Unwinding, $F=\Psi\circ h$ with $\Psi(\zeta)=\varphi_{-w}\bigl(\varphi_{-b}(\zeta)^2\bigr)$, a holomorphic map $D\to D$ with $\Psi(0)=F(z_0)=0$. $\Psi$ is not injective, because squaring is not, so it is not a rotation,
and the Schwarz lemma gives $|\Psi'(0)|<1$. Then $|F'(z_0)|=|\Psi'(0)||h'(z_0)|<|h'(z_0)|$, contradicting maximality.

\emph{Normalisation and uniqueness.} Multiplying $F$ by $e^{-i\arg F'(z_0)}$ makes $F'(z_0)>0$, and $F^{-1}$ is holomorphic by Lemma~\ref{lem:injective}. If $G$ is another such map, $\psi=F\circ G^{-1}$ maps $D$
onto $D$ with $\psi(0)=0$; the Schwarz lemma applied to $\psi$ and $\psi^{-1}$ gives $|\psi'(0)|=1$, so $\psi(z)=cz$, and $c=\psi'(0)=\frac{F'(z_0)}{G'(z_0)}>0$ forces $c=1$.
\end{proof}

\begin{corollary}\label{cor:jordanconformal}
The inside of every Jordan curve is mapped onto the unit disc by a holomorphic bijection with holomorphic inverse.
\end{corollary}
\begin{proof}
Volume~7 proves that the inside of a Jordan curve is connected, open and bounded, hence not $\C$. Volume~8 proves that this inside is homologically simply connected, so the theorem applies.
\end{proof}



\begin{think}
The proof of Theorem~\ref{thm:riemannmap} is an existence proof: it finds $F$ by maximising $|f'(z_0)|$, with no formula. Why does maximising the derivative at one point force the map to be onto?
(Picture what $|f'(z_0)|$ measures, and what happens to it when the image of $f$ is a disc with a hole.)
\end{think}


\section{Exercises}

\begin{exercise}[$\star$]\label{ex:riemannmaps}
\emph{Explicit Riemann maps.} Find holomorphic bijections onto the unit disc $D$ from: (a) the upper half-plane (try $z\mapsto\frac{z-i}{z+i}$); (b) the strip $|\operatorname{Im}z|<\frac\pi2$; (c) the slit plane
$\C\setminus\left]-\infty,0\right]$. Check that your maps are bijective.
\end{exercise}

\section{Hints}

\begin{hint}{ex:riemannmaps}
(a) Compare the distances from $z$ to $i$ and to $-i$. (b) The exponential straightens the strip into a half-plane. (c) A square root opens the slit plane into a half-plane.
\end{hint}

\section{Solutions}

\begin{solution}{ex:riemannmaps}
\idea Build the map from simple pieces that each straighten one feature: exponentials straighten strips, roots open angles, and a M\"obius map bends a half-plane into a disc.
\writeup (a) $|z-i|<|z+i|$ exactly when $z$ is closer to $i$ than to $-i$, that is, $\operatorname{Im}z>0$; so $z\mapsto\frac{z-i}{z+i}$ maps the upper half-plane $\mathbb H$ into $D$. Its inverse $w\mapsto i\frac{1+w}{1-w}$
maps $D$ into $\mathbb H$, since $\operatorname{Im}\bigl(i\frac{1+w}{1-w}\bigr)=\operatorname{Re}\frac{1+w}{1-w}=\frac{1-|w|^2}{|1-w|^2}>0$. (b) $e^z$ maps the strip bijectively onto the right half-plane ($e^{x+iy}$ has modulus $e^x$ and argument
$y\in\left]-\frac\pi2,\frac\pi2\right[$); multiplying by $i$ gives $\mathbb H$, and (a) gives $\frac{ie^z-i}{ie^z+i}=\frac{e^z-1}{e^z+1}=\tanh\frac z2$. (c) The principal square root maps the slit plane bijectively onto the right
half-plane, and $w\mapsto\frac{w-1}{w+1}$ maps the right half-plane onto $D$ (as in (a), after a rotation); so $z\mapsto\frac{\sqrt z-1}{\sqrt z+1}$.
\end{solution}

\section*{Sources and further reading}
\addcontentsline{toc}{section}{Sources and further reading}

B.~Riemann, \emph{Grundlagen f\"ur eine allgemeine Theorie der Functionen einer ver\"anderlichen complexen Gr\"osse} (G\"ottingen, 1851); W.~F.~Osgood, ``On the existence of the Green's
function for the most general simply connected plane region'', \emph{Trans.\ Amer.\ Math.\ Soc.} 1 (1900), 310--314; C.~Carath\'eodory, ``Untersuchungen \"uber die konformen Abbildungen von
festen und ver\"anderlichen Gebieten'', \emph{Math.\ Ann.} 72 (1912), 107--144; P.~Montel, ``Sur les suites infinies de fonctions'', \emph{Ann.\
Sci.\ \'Ec.\ Norm.\ Sup.} 24 (1907), 233--334; T.~Rad\'o, \"Uber die Fundamentalabbildung schlichter Gebiete'', \emph{Acta Sci.\ Math.\ (Szeged)} 1 (1923), 240--251. The history of the mapping theorem, with Koebe's and Fej\'er--Riesz's contributions, is in R.~Remmert,
\emph{Classical Topics in Complex Function Theory} (Springer, 1998), Chapter~8.
